\section{Implementation}
\label{sec:implementation}

\subsection{Crate Architecture}

The \re{} crate is structured as three source files within the
\texttt{redist-ensemble} module of the \texttt{redist} Rust workspace:

\begin{description}
  \item[\texttt{spanning.rs}] Wilson's loop-erased random walk.
    Exports \texttt{wilson\_ust(graph, rng)} returning an adjacency-list
    spanning tree. Uses \texttt{SmallRng} for per-step randomness (see
    Section~\ref{ssec:rng}).

  \item[\texttt{recom.rs}] The \recom{} proposal.
    Exports \texttt{recom\_step(partition, graph, pop, tol, rng)},
    which selects a pair, calls \texttt{wilson\_ust}, enumerates balanced
    cuts, and returns the updated partition.
    Handles tree-resample-on-failure and pair-reselection logic.

  \item[\texttt{chain.rs}] Parallel runner.
    Exports \texttt{run\_ensemble(config)} which drives $C$ independent
    chains via Rayon's \texttt{par\_iter}, collecting per-step statistics
    and computing convergence diagnostics.
\end{description}

\subsection{Graph Representation}

The census-tract adjacency graph is stored as a \texttt{Vec<Vec<usize>>}
adjacency list, loaded from the \texttt{.adj.bin} binary format shared
with \texttt{redist-data}.
Population data is a \texttt{Vec<u64>} aligned to the same vertex
indices.
Subgraphs for individual district pairs are constructed on demand as
\texttt{Vec<Vec<usize>>} slices by remapping vertex indices to the
compact range $\{0,\ldots,m-1\}$.
This subgraph construction cost is $O(m)$ per step and does not allocate
on the heap for graphs with $m \leq 512$ due to a stack-based small-buffer
optimization.

\subsection{Random Number Generation}
\label{ssec:rng}

Each chain uses a \texttt{SmallRng} instance from the \texttt{rand} crate
(version 0.8), seeded from the SHA-256-derived per-chain seed as described
in Section~\ref{sec:algorithm}.
\texttt{SmallRng} is a 128-bit xoshiro128++ generator: its state is small
(16 bytes), its throughput is high ($\sim$4 billion integers per second
on modern hardware), and it passes the PractRand statistical test suite
up to large byte counts (BigCrush status for the specific xoshiro128++
variant should be verified against \texttt{rand} crate documentation;
we make no claim here).
For the random walk within Wilson's algorithm, the dominant operation is
selecting a uniformly random neighbor of the current vertex.
We avoid modulo bias by using Lemire's fast divisor trick~\citep{lemire2019fast}
as implemented in \texttt{rand::Rng::gen\_range}.

\subsection{Balance Check}

The balance check for a candidate spanning tree $T$ of the merged region
is implemented as a single depth-first traversal.
We root $T$ at vertex 0 of the local indexing and compute subtree
population sums in $O(m)$ time using a post-order traversal.
For each edge $(u, \mathit{parent}(u))$, the subtree rooted at $u$ has
population $\mathit{pop}[u]$, and the complementary subtree has
population $P_R - \mathit{pop}[u]$, where $P_R$ is the total population
of the merged region.
The edge is \emph{balanced} if both subtrees satisfy
$|\mathit{pop}/\bar{p} - 1| \leq \varepsilon$.
Balanced edges are collected into a \texttt{SmallVec<[usize; 8]>} to
avoid heap allocation for the common case where the number of balanced
cuts is small.

\subsection{Serde Output}

Step-level statistics---district population variances, normalized cut
fraction, Democratic seat share (if election data is loaded)---are
serialized to newline-delimited JSON using \texttt{serde\_json}.
Chain-level diagnostics ($\rhat$, $\ess$, $\ham(1)$) are written at
chain completion.
The output schema is versioned: the constant
\texttt{ensemble\_output\_version: "1.0"} is emitted as a header field,
enabling downstream consumers to detect format changes.
The output format is designed to be streamed directly into the
\texttt{redist-report} pipeline for court-submission documents.

\subsection{Texas and Bipartition Failures}
\label{ssec:tx-impl}

Texas ($n = 4{,}866$ tracts, $k = 38$ districts) presents a known
challenge: the state's shape and high $k$ mean that a randomly merged
pair of adjacent districts often contains a region with no balanced
bipartition of any random spanning tree.
This is a combinatorial property of the Texas plan space, not a Python
limitation.
\gc{} without pair reselection is documented to fail on Texas from a cold
start~\citep{gerrychain2021}; \gc{} \emph{with} pair reselection would
also handle these cases, albeit more slowly.

\re{} handles this via the pair-reselection mechanism in
Algorithm~\ref{alg:recom}: after 10 consecutive spanning tree resample
failures for a given pair, a new adjacent pair is selected.
In practice, for Texas, pair reselection triggers frequently during the
first 50--100 steps from a cold start, then subsides as the chain moves
away from the starting plan into regions of higher feasibility density.
\re{}'s advantage over Python \gc{} is throughput: pair reselection at
microsecond speed (rather than second speed) reduces warm-up wall time
without altering the algorithm's combinatorial behavior.
The 10-resample budget is exhausted in microseconds in Rust rather than
the seconds it would take in Python, so pair reselection adds negligible
wall-time overhead.

The stationarity implications of pair reselection are discussed in
Section~\ref{sec:algorithm}.
